\section{Related Work}
\label{sec:related_work}

\subsection{Vision-based Social Navigation}
Navigating socially compliant robots in human-populated environments is a well-established challenge. Traditional approaches and benchmarks focus on avoiding collisions while respecting social norms using visual observations \cite{SEANSocialEnvironment2020, SocNavBenchGroundedSimulation2022, IGibsonChallenge2021}. Recent advances employ predictive models to forecast human trajectories. For instance, Falcon \cite{CognitionPrecognitionFutureaware2025} utilizes future-aware forecasting to anticipate human motion, while other works leverage Large Language Models (LLMs) to reason about social context and human intent \cite{CoNavBenchmarkHumancentered2024, OpenNavOpenworldNavigation2025}. 
However, a critical limitation of these vision-centric methods is their reliance on Line-of-Sight (LOS) perception. They typically fail to account for risks that are temporarily occluded, leading to unsafe interactions in ``blind corner'' scenarios. In contrast, SAVNav integrates audio perception to transcend LOS limitations, allowing the robot to anticipate unseen agents through their acoustic footprint before they become visible.

\subsection{Audio-Visual Navigation and Mapping}
The introduction of high-fidelity acoustic simulators, such as SoundSpaces 2.0 \cite{SoundSpaces20Simulation2022} and SonicSim \cite{SonicSimCustomizableSimulation2024}, has significantly accelerated research in Audio-Visual Navigation (AV-Nav). The standard paradigm involves agents learning to locate sounding targets (e.g., ringing phones) using reinforcement learning \cite{LookListenAct2020, AudiovisualNavigationNoisy2025}. Parallel efforts focus on constructing multi-modal maps. Methods like AVLMaps \cite{AudioVisualLanguage2024} and ConceptFusion \cite{ConceptFusionOpensetMultimodal2023} embed audio-visual features into 3D representations to support semantic queries, often leveraging foundation models like CLIP \cite{LearningTransferableVisual2021} or ImageBind \cite{ImageBindOneEmbedding2023}.
Despite these advances, existing AV-Nav methods primarily treat sound as a goal beacon (navigation target) or a static semantic attribute. They largely overlook the potential of sound as a dynamic collision warning signal. Furthermore, map-based approaches often struggle with highly dynamic entities \cite{Multimodal3DMap2024}. Our work diverges by treating audio as a predictor of dynamic occupancy. Instead of merely finding a sound source, we introduce \textit{Acoustic Hallucinated Momentum} to proactively project risk fields for collision avoidance.

\subsection{Active Perception and Multi-modal Reasoning}
To operate robustly, robots must handle perceptual uncertainty. Multi-modal Large Language Models (MLLMs) like Gemini \cite{Gemini25Pushing2025} and Qwen-Omni \cite{Qwen3omniTechnicalReport2025} demonstrate strong reasoning capabilities across audio and vision. Recent work such as BAT \cite{BATLearningReason2024} explores reasoning about spatial environments directly from acoustic echoes and sounds. Similarly, RILA \cite{RILAReflectiveImaginative2024} investigates reflective and imaginative search strategies for object localization.
While these works enhance semantic understanding, they often lack a mechanism for real-time kinematic reaction to uncertainty. In social navigation, the ambiguity of a sound source's direction is not just an informational gap but a safety risk. SAVNav addresses this by incorporating an Information-Theoretic Active Search strategy that explicitly drives the robot to viewpoints that maximize information gain regarding dynamic acoustic threats, balancing exploration with social safety.
